\documentclass[11pt]{article}
\usepackage{amsmath,amssymb,amsthm,mathtools}
\usepackage[margin=1in]{geometry}
\usepackage{enumitem}
\usepackage{hyperref}

\newtheorem{theorem}{Theorem}[section]
\newtheorem{lemma}[theorem]{Lemma}
\newtheorem{proposition}[theorem]{Proposition}
\newtheorem{corollary}[theorem]{Corollary}
\theoremstyle{definition}
\newtheorem{definition}[theorem]{Definition}
\theoremstyle{remark}
\newtheorem{remark}[theorem]{Remark}

\newcommand{\Hb}{\mathbb{H}}
\newcommand{\bp}{\boxplus}
\newcommand{\bigboxplus}{\mathop{\mbox{\normalsize$\boxplus$}}\limits}

\title{Problem 5 -- Humanize}
\author{}
\date{}

\begin{document}
\maketitle

\begin{abstract}
\noindent\textbf{Problem.} Let $\Hb$ be a hyperfield with hyperaddition
$\boxplus$. For $x,y,z\in\Hb$ with $0\in x\bp y\bp z$, does it follow that
$0\in x^{3}\bp y^{3}\bp z^{3}\bp(-xyz)\bp(-xyz)\bp(-xyz)$?

\smallskip\noindent\textbf{Content of this document.} We give complete
proofs, from the hyperfield axioms alone, of the degenerate case $xyz=0$
(Proposition~\ref{prop:degen}) and of the case in which one of the three
pair sums $x\bp y$, $x\bp z$, $y\bp z$ is a singleton
(Corollary~\ref{cor:collapse}). The remaining case is isolated in the
precisely stated Elevation Problem~\ref{prob:elevation}, which is
equivalent to the unrestricted statement; the document proves the
conditional implication (Elevation Problem $\Rightarrow$ full theorem) in
Theorem~\ref{thm:main}. We do not assert the unrestricted answer ``yes''
as proved in this round; no computation, and no external theorem
or unproved lemma is used as a proof step anywhere below.
\end{abstract}

\section{Preliminaries}\label{sec:prelim}

\begin{definition}\label{def:hf}
A \emph{hyperfield} is a set $\Hb$ with an element $0$, a commutative
multiplication making $\Hb^\times:=\Hb\setminus\{0\}$ an abelian group
(with $0\cdot a=0$), the abbreviation $-a:=(-1)a$, and a hyperaddition
$\bp:\Hb\times\Hb\to\mathcal P(\Hb)\setminus\{\varnothing\}$ satisfying:
\begin{enumerate}[label=(H\arabic*)]
\item $x\bp y=y\bp x$;\; $x\bp0=\{x\}$;\; $0\in x\bp y\iff y=-x$;
\item $(x\bp y)\bp z=x\bp(y\bp z)$, where for $S\subseteq\Hb$ one sets
$S\bp z:=\bigcup_{s\in S}s\bp z$. All bracketings of an iterated sum
coincide; we write $x_1\bp\cdots\bp x_n$, and $a^{\bp k}$ for the $k$-fold
hypersum of $a$ with itself;
\item $a(x\bp y)\subseteq ax\bp ay$;
\item $t\in x\bp y\implies x\in t\bp(-y)$.
\end{enumerate}
\end{definition}

\begin{lemma}[elementary consequences]\label{lem:basic}
Let $a,x,y,z\in\Hb$, $a\neq0$.
\begin{enumerate}[label=\textup{(\roman*)}]
\item $a(x\bp y)=ax\bp ay$;
\item $-(x\bp y)=(-x)\bp(-y)$;
\item $0\in x\bp y\bp z\iff -z\in x\bp y$;
\item if $x\bp y=\{t\}$ is a singleton with $x,y\neq0$, then
$-y\in x\bp(-t)$ and $-x\in y\bp(-t)$;
\item if $0\in x\bp y\bp z$ then $0\in(ax)\bp(ay)\bp(az)$.
\end{enumerate}
\end{lemma}

\begin{proof}
(i) (H3) gives $a(x\bp y)\subseteq ax\bp ay$; applying (H3) with scalar
$a^{-1}$ gives $a^{-1}(ax\bp ay)\subseteq x\bp y$, hence the reverse
inclusion $ax\bp ay\subseteq a(x\bp y)$.
(ii) is (i) with $a=-1$.
(iii) If $0\in x\bp y\bp z$, pick $t\in x\bp y$ with $0\in t\bp z$; by
(H1), $z=-t\in-(x\bp y)$, so $-z\in x\bp y$ by (ii). Conversely, if
$-z\in x\bp y$ then $0\in(-z)\bp z\subseteq x\bp y\bp z$ by (H1).
(iv) We have $t\in x\bp y$. By (H4), $x\in t\bp(-y)$; then
$0\in x\bp(-x)\subseteq(t\bp(-y))\bp(-x)$ by (H1). The hypersum
$(t\bp(-y))\bp(-x)$ contains $0$, so by (iii)
\[
  -t\ \in\ (-y)\bp(-x)\ =\ -(x\bp y)\ =\ \{-t\},
\]
the middle equality being (ii). Hence $-t\in(-x)\bp(-y)$; applying (H4)
to this membership gives $-y\in x\bp(-t)$. The proof of
$-x\in y\bp(-t)$ is symmetric, exchanging the roles of $x$ and $y$.
(v) From $-z\in x\bp y$ (iii) and (i) we get $-az\in ax\bp ay$, and
(iii) again gives the claim.
\end{proof}

For a finite multiset $M$ we write $\boxplus M$ for the (unambiguous)
hypersum of its elements.

\begin{lemma}[substitution atoms]\label{lem:atoms}
Let $a,b,c\in\Hb^{\times}$ with $-c\in a\bp b$.
\begin{enumerate}[label=\textup{(A\arabic*)}]
\item $0\in a\bp b\bp c$ and, for every $k\ge1$,
$0\in a^{\bp k}\bp b^{\bp k}\bp c^{\bp k}$.
\item \textup{(cubing a pair sum)} With
$(a\bp b)^{\cdot3}:=\{t_1t_2t_3:t_i\in a\bp b\}$,
\[
  -c^{3}\ \in\ (a\bp b)^{\cdot3}\ \subseteq\
   a^{3}\bp\bigl(a^{2}b\bigr)^{\bp3}\bp\bigl(ab^{2}\bigr)^{\bp3}\bp b^{3}.
\]
\item For every $\lambda\in\Hb^\times$, $-\lambda c\in\lambda a\bp\lambda b$.
\end{enumerate}
\end{lemma}

\begin{proof}
(A1) The case $k=1$ is Lemma~\ref{lem:basic}(iii). For general $k$, use the
same value $-c\in a\bp b$ in each of $k$ blocks: by (H2) the $3k$-term
hypersum obtained from $k$ disjoint blocks each isomorphic to
$a\bp b\bp c$ contains $0$ (realize each block by the value $-c$ and each
pair of block-outputs by $0$); by (H1) this multiset is exactly
$a^{\bp k}\cup b^{\bp k}\cup c^{\bp k}$.

(A2) Choosing $t_1=t_2=t_3=-c$ gives $-c^3\in(a\bp b)^{\cdot3}$. For the
inclusion, let $u,v\in a\bp b$. By Lemma~\ref{lem:basic}(i),
\[
  uv\ \in\ u(a\bp b)\ =\ ua\bp ub\ \subseteq\
   \bigl(a^2\bp ab\bigr)\bp\bigl(ab\bp b^2\bigr)
   \ =\ a^2\bp(ab)^{\bp2}\bp b^2 ,
\]
using $u\in a\bp b$ to get $ua\in a^2\bp ab$ and $ub\in ab\bp b^2$. Now
any $t=t_1t_2t_3\in(a\bp b)^{\cdot3}$ is $(t_1t_2)t_3$ with
$u:=t_1t_2\in a^2\bp(ab)^{\bp2}\bp b^2$ and $t_3\in a\bp b$, whence
\[
  t=ut_3\ \in\ \bigl(a^2\bp(ab)^{\bp2}\bp b^2\bigr)\,t_3
  \ \subseteq\ a^3\bp(a^2b)^{\bp3}\bp(ab^2)^{\bp3}\bp b^3,
\]
since $\gamma t_3\in\gamma a\bp\gamma b$ for each
$\gamma\in\{a^2,ab,b^2\}$ (Lemma~\ref{lem:basic}(i)), and (H1) collects
the exponents.

(A3) is Lemma~\ref{lem:basic}(i) applied to $-c\in a\bp b$.
\end{proof}

\section{Statement and degenerate case}\label{sec:easy}

\begin{theorem}[conditional form]\label{thm:main}
Let $\Hb$ be a hyperfield and $x,y,z\in\Hb$ with $0\in x\bp y\bp z$. Then
\[
  0\ \in\ \Sigma\ :=\
   x^3\bp y^3\bp z^3\bp(-w)\bp(-w)\bp(-w),\qquad w:=xyz,
\]
provided the Elevation Problem \ref{prob:elevation} below has an
affirmative answer. The degenerate instance $w=0$ and the collapse
instance (some pair sum among $\{x,y,z\}$ a singleton) are unconditional
(Proposition~\ref{prop:degen}, Corollary~\ref{cor:collapse}).
\end{theorem}

\begin{lemma}[scalar invariance]\label{lem:scale}
For $\lambda\in\Hb^\times$, the implication
$0\in x\bp y\bp z\Rightarrow0\in\Sigma$ holds for $(x,y,z)$ iff it holds
for $(\lambda x,\lambda y,\lambda z)$.
\end{lemma}
\begin{proof}
The hypothesis is invariant by Lemma~\ref{lem:basic}(v). The conclusion
hypersum of the scaled triple equals $\lambda^3\Sigma$ by repeated use of
Lemma~\ref{lem:basic}(i), and $0\in\lambda^3\Sigma\iff0\in\Sigma$.
\end{proof}

\begin{proposition}[degenerate case]\label{prop:degen}
If $xyz=0$ and $0\in x\bp y\bp z$, then $0\in\Sigma$.
\end{proposition}
\begin{proof}
Permuting the variables, we may assume $z=0$. For $t$ ranging over
$x\bp y$, (H1) gives $t\bp z=\{t\}$; since $0\in x\bp y\bp z$ some such
hypersum $t\bp z$ contains $0$, so $t=0$, and (H1) gives $y=-x$; then
$w=0$. If $x=0$ as well, all six terms of $\Sigma$ are $0$, and any
hypersum of zeros is $\{0\}$. Otherwise $y=-x\neq0$: then
$y^3=(-x)^3=-x^3$, and $z^3=0$, $-w=0$, so writing the hypersum in the
bracketing $\bigl(x^3\bp y^3\bigr)\bp z^3\bp(-w)\bp(-w)\bp(-w)$, the
first factor contains $0$ by (H1) (since $y^3=-x^3$), and appending the
remaining zero terms preserves the value $0$ (H1 again:
$0\bp0=\{0\}$). Hence $0\in\Sigma$.
\end{proof}

From now on $w=xyz\neq0$.

\section{The collapse case}\label{sec:collapse}

\begin{lemma}[collapse lemma]\label{lem:collapse}
Let $X,Y\in\Hb^\times$ with $X\bp Y=\{1\}$. Then
$0\in X^3\bp Y^3\bp(-1)\bp(XY)\bp(XY)\bp(XY)$.
\end{lemma}

\begin{proof}
Set $\mathcal U:=(X^2Y)^{\bp3}\bp(XY^2)^{\bp3}$ and
$\mathcal W:=(XY)^{\bp3}$. Scaling $X\bp Y=\{1\}$ by the scalar $XY$
(Lemma~\ref{lem:basic}(i)):
\[
  X^2Y\bp XY^2=\{XY\}.
\]
Regrouping the six terms of $\mathcal U$ into three such pairs (H1),(H2):
\begin{equation}\label{eq:star}
  \mathcal U
  \;=\;\bigboxplus_{i=1}^{3}\bigl(X^2Y\bp XY^2\bigr)
  \;=\;\{XY\}\bp\{XY\}\bp\{XY\}
  \;=\;\mathcal W .
\end{equation}
Apply (A2) to $1\in X\bp Y$ (so $c=-1$ and $-c^3=1$):
\[
  1\ \in\ (X\bp Y)^{\cdot3}\ \subseteq\
   X^3\bp(X^2Y)^{\bp3}\bp(XY^2)^{\bp3}\bp Y^3 ,
\]
so there exist $g\in X^3\bp Y^3$ and $u\in\mathcal U$ with $1\in g\bp u$;
by (H1), $0\in g\bp u\bp(-1)$, i.e.\
\begin{equation}\label{eq:i1}
  0\ \in\ g\bp u\bp(-1)\qquad
   \text{for some } g\in X^3\bp Y^3,\ u\in\mathcal U .
\end{equation}
By \eqref{eq:star} the value $u$ is an element of
$\mathcal W=(XY)^{\bp3}$, so the singleton $\{u\}$ is included in the
value set of the three-term hypersum $(XY)\bp(XY)\bp(XY)$; value
substitution (H2) in \eqref{eq:i1} therefore gives
\[
  0\ \in\ g\bp(-1)\bp(XY)\bp(XY)\bp(XY)
  \ \subseteq\ X^3\bp Y^3\bp(-1)\bp(XY)^{\bp3} . \qedhere
\]
\end{proof}

\begin{corollary}[collapse case of Theorem~\ref{thm:main}]\label{cor:collapse}
Let $0\in x\bp y\bp z$ with $w=xyz\neq0$, and suppose some pair sum among
$\{x,y,z\}$ is a singleton. Then $0\in\Sigma$.
\end{corollary}
\begin{proof}
Permuting the variables, assume $x\bp y$ is a singleton. Since
$-z\in x\bp y$ (Lemma~\ref{lem:basic}(iii)), that singleton is $\{-z\}$.
Rescale by $\lambda:=-z^{-1}$ (Lemma~\ref{lem:scale}) and put
$X:=-z^{-1}x$, $Y:=-z^{-1}y$: the scaled triple is $(X,Y,-1)$ with
$X\bp Y=\{1\}$ (Lemma~\ref{lem:basic}(i)), and the scaled value of $w$ is
$(-z^{-1})xyz=-XY$. By Lemma~\ref{lem:scale} the conclusion for the
scaled triple, namely
$0\in X^3\bp Y^3\bp(-1)\bp(XY)\bp(XY)\bp(XY)$, is exactly
Lemma~\ref{lem:collapse}, and it pulls back to $0\in\Sigma$ by
Lemma~\ref{lem:scale} once more.
\end{proof}

\section{The free case: towers and the elevation problem}\label{sec:free}

Assume throughout:
\begin{equation}\label{eq:free}
  0\in x\bp y\bp z,\qquad w=xyz\neq0,\qquad
  x\bp y\neq\{-z\},\;x\bp z\neq\{-y\},\;y\bp z\neq\{-x\}.
  \tag{F}
\end{equation}
Set
\[
  P:=x^2y\bp xy^2,\qquad H:=x^3\bp y^3\bp z^3,\qquad
  \overline V:=(-w)\bp(-w)\bp(-w).
\]
Scaling $-z\in x\bp y$ by $xy$ (Lemma~\ref{lem:basic}(i)) gives
$-w\in P$. The goal is $0\in H\bp\overline V$.

\begin{lemma}[the two towers]\label{lem:towers}
Under \eqref{eq:free}:
\begin{enumerate}[label=\textup{(\roman*)}]
\item $0\in H\bp P\bp P\bp P$, i.e.\ there exist $h\in H$ and
$p_1,p_2,p_3\in P$ with $0\in h\bp p_1\bp p_2\bp p_3$;
\item $0\in P\bp P\bp P\bp w\bp w\bp w$.
\end{enumerate}
\end{lemma}

\begin{proof}
(i) Apply (A2) to $-z\in x\bp y$:
$-z^3\in x^3\bp(x^2y)^{\bp3}\bp(xy^2)^{\bp3}\bp y^3$, whence by
Lemma~\ref{lem:basic}(iii)
$0\in H\bp(x^2y)^{\bp3}\bp(xy^2)^{\bp3}$ after regrouping (H1). The six
mixed terms regroup into three pairs (H1),(H2):
\begin{equation}\label{eq:blockP}
  (x^2y)^{\bp3}\bp(xy^2)^{\bp3}
  \;=\;\bigboxplus_{i=1}^{3}\bigl(x^2y\bp xy^2\bigr)
  \;=\;P\bp P\bp P ,
\end{equation}
which gives (i).
(ii) Apply (A1) with $k=3$ to $-w\in P$: the blocks are the relation
$-w\in x^2y\bp xy^2$, i.e.\ $0\in x^2y\bp xy^2\bp w$, three times; their
union regroups (H1) into $P\bp P\bp P\bp w^{\bp3}$.
\end{proof}

\begin{lemma}[free propagation]\label{lem:prop}
Under \eqref{eq:free}, each of the following pair sums has at least two
elements:
\[
  x^2y\bp xy^2,\qquad x^2z\bp xz^2,\qquad y^2z\bp yz^2,\qquad
  C\bp w,\ \ C\bp(-w)\ \ \text{for }\ 
  C\in\{x^2y,xy^2,x^2z,xz^2,y^2z,yz^2\}.
\]
\end{lemma}

\begin{proof}
Each displayed pair sum is, by Lemma~\ref{lem:basic}(i), a scalar multiple
of one of the level-1 pair sums $x\bp y$, $x\bp z$, $y\bp z$: with
$w=xyz$,
\begin{align*}
  x^2y\bp xy^2&=xy\,(x\bp y), &
  x^2z\bp xz^2&=xz\,(x\bp z), &
  y^2z\bp yz^2&=yz\,(y\bp z),\\
  x^2y\bp w&=xy\,(x\bp z), &
  xy^2\bp w&=xy\,(y\bp z), &
  x^2z\bp w&=xz\,(x\bp y),\\
  xz^2\bp w&=xz\,(y\bp z), &
  y^2z\bp w&=yz\,(x\bp y), &
  yz^2\bp w&=yz\,(x\bp z),
\end{align*}
and by Lemma~\ref{lem:basic}(ii) each $C\bp(-w)$ is the negative of the
corresponding $C\bp w$. Scalar multiples preserve the cardinality of a
pair sum. Finally, each of $x\bp y$, $x\bp z$, $y\bp z$ has at least two
elements: $-z\in x\bp y$, so if $x\bp y$ were a singleton it would equal
$\{-z\}$, excluded by \eqref{eq:free}; similarly
$-y\in x\bp z$ and $-x\in y\bp z$ exclude the other two.
\end{proof}

\begin{remark}[what the towers do and do not give]\label{rem:gap}
Lemma~\ref{lem:towers}(i) says there exist $h\in H$ and
$p_1,p_2,p_3\in P$ with $0\in h\bp p_1\bp p_2\bp p_3$, while the target
is the special case $p_1=p_2=p_3=-w$ (note $-w\in P$). Replacing the
three values $p_i$ by the \emph{particular} value $-w$ is not licensed by
any axiom: the substitution rules (H2),(H3) propagate \emph{superset}
inclusions only, and value-level replacements (via (H4)) require
singleton pair sums, which Lemma~\ref{lem:prop} rules out in the free
case. In every computed finite model the replacement turns out to be
possible (see Remark~\ref{rem:honest}(d)), but no proof from (H1)--(H4)
was found in this attempt; this is exactly the content of the following
problem.
\end{remark}

\begin{remark}[the elevation problem; equivalent form]\label{prob:elevation}
\textbf{Elevation Problem (open).} Is it true that for every hyperfield
$\Hb$ and all $a,b\in\Hb^\times$,
\[
  1\ \in\ a\bp b\ \Longrightarrow\
   0\ \in\ a^{3}\bp b^{3}\bp(-1)\bp(ab)\bp(ab)\bp(ab)\ ?
   \tag{E}
\]
\textbf{Equivalence with the free case.} (The six terms of the displayed
hypersum are the element \emph{cubes} $a^3,b^3$ and $ab$, not threefold
hypersums: the stronger variant with $a^{\bp3}\bp b^{\bp3}$ in place of
$a^3\bp b^3$ is \emph{false} already in fields, while the element-cube
form above is exactly what scalar invariance transports. See
Remark~\ref{rem:honest}.) Let $0\in x\bp y\bp z$ with
$w=xyz\neq0$ in the free case \eqref{eq:free}; put $a:=-z^{-1}x$,
$b:=-z^{-1}y$. Then $-z\in x\bp y$ gives $1\in a\bp b$ (scale by
$-z^{-1}$), and scalar invariance (Lemma~\ref{lem:scale} with
$\lambda=-z^{-1}$) identifies the target hypersum with that of (E):
indeed $\lambda x=a$, $\lambda y=b$, $\lambda z=-1$, and
$\lambda^3(-w)=-(\lambda^3w)$ with
$\lambda^3w=(-z^{-1})^3xyz=-z^{-3}xyz=-xy=-ab$, so $\lambda^3(-w)=ab$.
Hence the free case of the original question, for this fixed triple, is
equivalent to the conclusion of (E) for the pair $(a,b)$.
\end{remark}

\begin{proof}[Proof of Theorem~\ref{thm:main} given the Elevation Problem]
Degenerate and collapse cases are
Proposition~\ref{prop:degen} and Corollary~\ref{cor:collapse}. In the
free case \eqref{eq:free}, apply the equivalence of
Proposition~\ref{prob:elevation}: the affirmative answer to (E) yields
\[
  0\ \in\ X^{3}\bp Y^{3}\bp(-1)\bp(XY)\bp(XY)\bp(XY)
\]
for $(X,Y)=(-z^{-1}x,-z^{-1}y)$, and Lemma~\ref{lem:scale} pulls this
precisely back to $0\in\Sigma$.
\end{proof}

\section{Concluding assessment}\label{sec:concl}

\begin{remark}[the zero-trap; why the remaining case is hard]
\label{rem:zerotrap}
The analysis of Sections \ref{sec:free} isolates exactly one algebraic
phenomenon separating the proved cases from the Elevation Problem.
Write the free-case target, after rescaling, as
$0\in a^{3}\bp b^{3}\bp(-1)\bp v$ for a suitable
$v\in(ab)^{\bp3}$. The tower relations available from the axioms
(Lemma~\ref{lem:towers} and its refinements) create the desired zero as
a value of a long hypersum, but every such zero arises, in the last
pair-sum step, in one of two shapes:
\begin{enumerate}[label=\textup{(Z\arabic*)}]
\item \emph{paired zero}: a value $t$ with $-t\in u$ for a known pair
  sum $u$, so that $0\in t\bp u$; these zeros are exactly the ones the
  substitution calculus can move around, and in all solved cases the
  witness is of this shape;
\item \emph{unpaired zero}: an element $x$ with $0\in x\bp x$ but
  $-x\notin x\bp x$ (note that $x\bp x$ may contain further elements).
  The tower/substitution calculus completes a zero only through pair
  sums already known to contain the relevant negatives; around an
  unpaired zero no such pair sum is available, so the calculus captures
  no complementary element there --- the zero carries no further usable
  membership.
\end{enumerate}
For example, in the quotient hyperfield $\mathbb F_{19}/\{\pm1\}$ the
element $x=\{8,11\}$ satisfies $x\bp x=\{0,\{3,16\}\}$, so
$0\in x\bp x$ while $x$ (which equals $-x$) itself is not in
$x\bp x$; and in the free instance $1\in a\bp b$ with $a=\{2,17\}$,
$b=\{3,16\}$ (for which $a^{3}=b^{3}=\{8,11\}$) the completing zero of
the Elevation Problem is realized at the value $\{1,18\}$, which is
self-inverse and unpaired ($0\in\{1,18\}\bp\{1,18\}$ but
$\{1,18\}\notin\{1,18\}\bp\{1,18\}$), so no realization through a
paired zero exists for this instance. Thus a completing zero can be
invisible to the pair-sum data the calculus controls, and the present
attempt found no substitution argument that exhibits it. Proving the
Elevation Problem therefore appears to require either a genuinely
non-substitution argument or a proof that, whenever $1\in a\bp b$, one
can always force a paired realization of $0\in
a^{3}\bp b^{3}\bp(-1)\bp(ab)^{\bp3}$; neither was achieved here.
\end{remark}

\begin{remark}[status and atlas of failed routes]\label{rem:honest}
(a) Everything proved in this document follows from (H1)--(H4) alone; no
external theorem, no computation, and no unproved lemma is used anywhere.

(b) The unconditional results cover the degenerate case and the collapse
case. The remaining free case is equivalent to the corrected (element-cube)
Elevation Problem~\ref{prob:elevation}, whose every symbol is defined in
Section~\ref{sec:prelim}. This problem marks exactly what is and is not
proved here; the document does \emph{not} claim the answer to the
original question in full (that requires solving the Elevation Problem).

(c) Review history. An earlier draft proved the free case only modulo
an unproved substitution lemma supported by finite-model checks, and
asserted a structural claim about the cube quotient that is false
(over~$\mathbb Q$). Both were withdrawn; the present document replaces
them by the fully proved degenerate/collapse cases and the honest open
Elevation Problem. The route atlas below records what the search
established.

(d) For reference, the question was also verified on all computed
models --- the quotient hyperfields $\mathbb F_{p}/U$
($271\,139$ zero-relations) and the small hyperfields of orders
$\le7$ ($618\,867$ instances in total), with no failure. This is
\emph{evidence, recorded for orientation only, not a proof step}: no
statement below depends on it.

(e) Failed routes toward the Elevation Problem, with the precise
obstruction found for each (proofs of the refuting statements are
routine verifications inside the named models and are left to the
reader; they are \emph{not} used as proof steps):
\begin{enumerate}[label=\textup{(\roman*)}]
  \item \emph{$1\boxplus1$ saturation}: the needed membership $1\in
  1\bp1$ fails: in $\mathbb F_{p}/U$ one has $1\in1\bp1$ iff $1+u\in U$
  for some $u\in U$, which fails for many~$U$ (30 of the 55 computed
  nontrivial models). Even when $1\in1\bp1$ holds, the derived sublemma
  ``$u\in1\bp1\Rightarrow1\in u\bp(-1)\bp(1\bp1)$'' fails in some of
  them.
  \item \emph{permutation-adjusted tower}: the permuted
  binomial regrouping produces the mixed block
  $(c_{1}c_{2})^{\bp3}\bp(c_{2}c_{3})^{\bp3}$, which splits into the
  null pairs $\{t,-t\}$ ($t=c_{1}c_{2}c_{3}$); no usable new relation
  remains.
  \item \emph{H4-collapse}: the conditional ``$c\bp d=\{1\}$, $t\in
  c\bp s$ $\Rightarrow$ $d\in s\bp(-t)$'' is false: in $\mathbb F_{5}$,
  take $c=2,d=4,s=1,t=3$: then $c\bp d=\{1\}$ and $t\in c\bp s=\{3\}$,
  but $d=4\notin s\bp(-t)=\{3\}$. The companion ``delta condition'' $-g\bp
  z^{\bp3}\subseteq a^{\bp3}\bp b^{\bp3}$ fails likewise.
   \item \emph{value-swap}: cubing $1\in a\bp b$ (atom (A2)) always
   yields some $g\in a^{3}\bp b^{3}$ and $m\in
   (a^{2}b)^{\bp3}\bp(ab^{2})^{\bp3}$ with $1\in g\bp m$, but the
   single-witness step $0\in g\bp(-1)\bp(ab)^{\bp3}$ fails for about a
   positive fraction of these ($g,m$) in the computed models, so this route does not
   close the Elevation Problem.
  \item \emph{substitution calculus}: a bounded forward search from the
  axiom-generated $0$-hypersums (closure under scaling, appending of
  null pairs $\{t,-t\}$, and substitution of proved pair-sum
  memberships) does not derive the target $0\in
  a^{3}\bp b^{3}\bp(-1)\bp(ab)^{\bp3}$ in checked models; this is a
  bounded search, not an impossibility theorem.
\end{enumerate}

(f) As to the expected answer: in a field the assertion holds by
$X^3+Y^3+Z^3-3XYZ=(X+Y+Z)(X^2+Y^2+Z^2-XY-YZ-ZX)$, and the hyperfield
statement is its multivalued analogue; the author's judgement (not a
proof claim; see (a)) is that it is true in every hyperfield. A proof
would follow from any solution of the Elevation Problem;
Remark~\ref{rem:zerotrap} records the algebraic feature (unpaired
zeros) that such a solution must control, which the substitution
calculus cannot supply.
\end{remark}

\end{document}
